% Ad Familiares 13.78 --- headnote
% ad-familiares-13-78, 46 BC, Rome

\textit{Cicero to Aulus Allienus, written from Rome at the
beginning of 46 BC (the manuscript dateline:
\textit{Scr. Romae in. a. 708 (46)}). Allienus was an
old associate of Cicero's; he succeeded Manius Acilius
Glabrio as proconsul of Sicily early in 45 BC, and this
letter looks toward that command, so although Perseus
prints an early-46 date the occasion belongs to the
threshold between Acilius's tenure and Allienus's. The
piece thus stands in succession to the Acilius cluster
of Fam.\ 13.30--39, on which Cicero had drawn for the
same province a few months before.}

\textit{The beneficiary is Democritus of Sicyon, named in
the opening sentence with a notice on his standing as a
Greek host: a rare degree of personal closeness on
Cicero's part to a man of that nation, and a tier-marker
within Cicero's own register of recommendations
(\textit{quod non multis contigit, Graecis praesertim}).
He is called nearly the foremost man of Achaea ---
\textit{paene Achaiae principem} --- a regional epithet
that places him at the top rank of the province. The
formal opening of the second section, ``I merely open the
way and clear the approach to your acquaintance''
(\textit{aditum ad tuam cognitionem patefacio et munio}),
is one of Cicero's recurring metaphors for the
introductory commendation; the close, with its
sequence ``embrace him, hold him dear, count him among
your own,'' moves the request into the warmest tier
of the genre.}
